%============================================================
% HALAMAN ABSTRAK
%============================================================
\cleardoublepage
\chapter*{ABSTRAK}
\addcontentsline{toc}{chapter}{ABSTRAK}
\thispagestyle{fancy}

Program Jaminan Kesehatan Nasional (JKN) menggunakan mekanisme pembayaran Indonesian Case-Based Groups (INA-CBG) pada fasilitas kesehatan rujukan tingkat lanjutan, sehingga nilai klaim sangat dipengaruhi oleh diagnosis, prosedur, tingkat keparahan, dan hasil grouping. Kondisi ini membuka risiko fraud atau upcoding yang dapat menimbulkan kerugian finansial dan memperberat proses verifikasi klaim. Permasalahan utama pada domain ini adalah sulitnya mengidentifikasi klaim berisiko secara cepat, terukur, dan dapat dijelaskan, terutama ketika label fraud aktual pada level klaim tidak tersedia.

Penelitian ini mengajukan pendekatan deteksi indikasi risiko fraud/upcoding berbasis machine learning dengan pendekatan claim-centric, yaitu menjadikan satu klaim sebagai unit analisis utama dan memperkaya klaim dengan informasi administratif, klinis, biaya, diagnosis sekunder, dan histori layanan. Untuk mengatasi keterbatasan label fraud aktual, penelitian ini membentuk target awal melalui pseudo-labeling berbasis Isolation Forest, kemudian membandingkan beberapa model supervised learning melalui benchmarking multi-model. Pipeline penelitian juga mencakup feature engineering, branch-based feature selection, threshold tuning, serta interpretabilitas global dan lokal.

Eksperimen dilakukan pada 1.699.219 klaim FKRTL yang telah diintegrasikan ke dalam dataset claim-centric. Hasil eksperimen menunjukkan bahwa konfigurasi terbaik diperoleh pada model XGBoost dengan branch VIF dan threshold 0,80, dengan ROC-AUC 0,999876, PR-AUC 0,996046, F1-score 0,952511, precision 0,913917, recall 0,994507, dan MCC 0,951883 terhadap target pseudo-label. Perlu dicatat bahwa seluruh metrik evaluasi pada penelitian ini dihitung terhadap pseudo-label sebagai target proksi, sehingga hasilnya tidak boleh ditafsirkan sebagai akurasi deteksi fraud aktual yang telah diverifikasi. Hasil interpretabilitas menunjukkan bahwa fitur biaya, tarif grouping, histori nonkapitasi, special CMG, dan rasio tarif merupakan faktor penting yang mendorong skor risiko. Dengan demikian, sistem yang dibangun berpotensi digunakan sebagai alat bantu analitis untuk memprioritaskan pemeriksaan klaim berisiko tinggi, bukan sebagai penetapan fraud final.

\vspace{1cm}
\noindent\textbf{Kata Kunci:} Fraud Detection, Upcoding, INA-CBG, Machine Learning, Claim-Centric
